\documentclass[10pt]{article}
\usepackage[margin=20mm]{geometry}
\usepackage{amsmath,amssymb,mathtools,fontspec,unicode-math,xurl,hyperref}
\IfFontExistsTF{DejaVu Serif}{\setmainfont{DejaVu Serif}}{\setmainfont{Latin Modern Roman}}
\setmathfont{latinmodern-math.otf}
\setlength{\parindent}{0pt}\setlength{\parskip}{6pt}\sloppy
\title{Technical Review of the Moment-Map Criterion and a Linear Cotangent Case}
\author{AI-assisted technical review}
\date{}
\begin{document}
\maketitle
Abstract: The equivalence in Proposition 2.1.1(1) fails with its stated allowance of an arbitrary locally closed invariant subset of the moment-map image. The standard Hamiltonian action of \(\operatorname{SL}_2\) on a symplectic plane gives a dimension-zero fiber satisfying the numerical condition while its restricted inverse image is not Lagrangian. The proof incorrectly treats a component of a restricted level as a component of the full Borel moment level. This is a substantive failure of the general criterion, not a counterexample to the paper’s case-by-case component counts. The ordinary linear cotangent argument in Section 4.1 passes the separate check described here. Independent acceptance and scientific adjudication remain pending.\par
\section*{Source and scope}
M25-0375 is Michael Finkelberg, Victor Ginzburg and Roman Travkin, Lagrangian Subvarieties of Hyperspherical Varieties, Geometric and Functional Analysis 35 (2025), 254–282, DOI 10.1007/s00039-025-00703-3 [source]. The assigned source is the 29-page version-of-record PDF; printed page \(p\) is physical PDF page \(p-253\).\par
This review checks the hypotheses and equivalences of Proposition 2.1.1, pp. 258–259, and the standard cotangent/rook argument for the linear side in Section 4.1, pp. 264–265. The exceptional and orthosymplectic enumerations, the complete maze bijections, and the equivariant-slice side are not independently certified. The broader statements in Section 1 remain distinguished as conjectures. In particular, the Weyl-group and categorical conjectures are not presented here as proved results.\par
\section*{Problem 1: restriction to a closed stratum destroys the claimed equivalence}
Classification: Major. Location: Proposition 2.1.1(1)(a)–(c), printed p. 258 (physical page 5), and the proof’s coisotropic-component assertion on p. 259 (physical page 6).\par
The proposition assumes that \(X\) is a smooth affine Hamiltonian \(G\)-variety, that \(Z\subset\mu(X)\) is a locally closed \(G\)-stable subset, and that a character \(\bar\lambda\in\mathfrak b^*\) has \(Z_{\bar\lambda}=Z\cap(\bar\lambda+\mathfrak b^\perp)\ne\varnothing\). It then declares the following equivalent: \[\begin{aligned}
 &\mu^{-1}(Z_{\bar\lambda})\text{ is Lagrangian in }X; \qquad\text{(a)}\\
 &\mu^{-1}(\lambda)\text{ is isotropic for every }\lambda\in Z_{\bar\lambda}; \qquad\text{(b)}\\
 &\dim F\leq \tfrac12\bigl(\dim X-\dim(G\lambda)\bigr)
 \quad\text{for every irreducible fiber component }F. \qquad\text{(c)}\end{aligned}\] There is no openness hypothesis on \(Z\) in part (1).\par
\subsection*{An explicit Hamiltonian counterexample}
Equation reference: eq:moment. Let \[G=\operatorname{SL}_2(\mathbb C),\qquad
 X=\mathbb A^2_{\mathbb C},\qquad
 \omega=dx\wedge dy,\] with the standard linear action, and let \(B\) be the upper-triangular Borel. Identify \(\mathfrak{sl}_2^*\) with \(\mathfrak{sl}_2\) by the trace pairing. A moment map is \[\mu(x,y)=\frac12
 \begin{pmatrix}xy&-x^2\\y^2&-xy\end{pmatrix}.
 \] For \(A=\left(\begin{smallmatrix}a&b\\c&-a\end{smallmatrix}\right)\), its component is \[\langle\mu,A\rangle=axy+\tfrac12by^2-\tfrac12cx^2,\qquad
 \iota_{A(x,y)}\omega=d\langle\mu,A\rangle .\] The matrix formula is \(\operatorname{SL}_2\)-equivariant. Equivalently, for the standard basis \(H,E,F\), the polynomial Hamiltonians \[\mu_H=xy,\quad \mu_E=y^2/2,\quad \mu_F=-x^2/2\] satisfy \(\{\mu_H,\mu_E\}=2\mu_E\), \(\{\mu_H,\mu_F\}=-2\mu_F\), and \(\{\mu_E,\mu_F\}=\mu_H\), with \(\{f,g\}=f_xg_y-f_yg_x\). Changing the overall moment-map sign convention would not affect any zero fiber below.\par
Choose \[Z=\{0\}\subset\mu(X),\qquad \bar\lambda=0.\] All hypotheses are satisfied: \(G\) is complex reductive, \(X\) is smooth affine symplectic, \(Z\) is closed and \(G\)-stable, and \(0\) is a character of \(\mathfrak b\). Moreover \(Z_{\bar\lambda}=\{0\}\). Equation [eq:moment] (\#eq:moment) gives \[\mu^{-1}(0)_{\mathrm{red}}=\{(0,0)\},\qquad
 \dim\mu^{-1}(0)=0,\qquad \dim(G\cdot0)=0.\] Consequently condition (c) is the true inequality \[0\leq\tfrac12(2-0)=1,\] whereas condition (a) is false: a point in a symplectic surface has dimension \(0\), not the required Lagrangian dimension \(1\). Thus (a) and (c) are not equivalent. This conclusion is independent of any convention about isotropy of a nonreduced scheme. With the usual reduced-variety interpretation, the point is also isotropic, so (b) holds as well.\par
The scheme fiber is defined by \((x^2,xy,y^2)\), but its sole irreducible component still has dimension \(0\); nilpotent structure cannot repair the stated dimension equivalence. The action is even coisotropic generically, since \(\operatorname{SL}_2\) is transitive on \(X\setminus\{0\}\). No failure of the general-orbit condition is being exploited.\par
\subsection*{The precise proof failure}
Equation reference: eq:borel. The full zero level of the \(B\)-moment map is \[\mu_B^{-1}(0)_{\mathrm{red}}
 =\{xy=0,\ y^2=0\}_{\mathrm{red}}
 =\{y=0\},
 \] which is a Lagrangian line. The restricted set \(\mu^{-1}(Z_{\bar\lambda})\) is the origin, a proper closed subset of that line and not an irreducible component of the full level.\par
The proof on p. 259 asserts that the component \(\Lambda\) under consideration is a component of the fiber of the \(B\)-moment map, and invokes coisotropy of that fiber component. That inference is invalid after restriction to an arbitrary locally closed \(Z\). In this example the origin is not coisotropic: its tangent space is zero and its symplectic orthogonal is the entire tangent plane. The preceding isotropy calculation therefore cannot be upgraded to the Lagrangian conclusion.\par
\subsection*{Impact and repair boundary}
The finding invalidates a central implication in the general moment-map criterion as actually quantified. It requires a restriction of the proposition’s scope or a different statement, rather than an arithmetic correction inside a valid proof; this is the basis for the Major classification.\par
An openness hypothesis on \(Z\), as already used in part (2), removes this particular obstruction. In the example, choosing \(Z=\mu(X)\setminus\{0\}\) gives \(\mu^{-1}(Z_0)=\{(x,0):x\ne0\}\), which is Lagrangian. More generally, openness makes \(\mu^{-1}(Z)\) open in \(X\), preventing the specific dimension collapse caused by selecting a smaller closed stratum. This is a repair direction, not an independent certification of every equivalence under a revised hypothesis.\par
The example does not disprove Proposition 2.1.1(2), the zero-\(B\)-level conjecture, or the numerical equalities in the later case studies. Indeed, its full zero-\(B\)-level [eq:borel] (\#eq:borel) is Lagrangian. The report does not infer that the paper’s component counts rely indispensably on the false arbitrary-\(Z\) extension.\par
\section*{Separate check: the ordinary linear cotangent side}
For \(V=\operatorname{Hom}(\mathbb C^N,\mathbb C^M)\) and the standard cotangent action, the defining identity for the moment map says \[\mu_B(v,\xi)=0
 \quad\Longleftrightarrow\quad
 \xi(\mathfrak b\cdot v)=0.\] Thus the zero level is the union of conormal bundles of the \(B_M\times B_N\)-orbits in \(V\). Each conormal bundle is isotropic and has dimension \(\dim O+\operatorname{codim}_V O=MN\). There are finitely many such orbits, indexed by partial permutation matrices, or non-attacking rook placements. Therefore their distinct closures are exactly Lagrangian components. This argument concerns the full Borel level and has no arbitrary closed-\(Z\) restriction.\par
The rook count is \[R_{M,N}=\sum_{r=0}^{\min(M,N)}
 \binom Mr\binom Nr r!.\] The southwest rectangular ranks are preserved by left and right upper-triangular row/column operations and determine a partial permutation matrix: each occupied entry is recovered by the corresponding second difference of the ranks. The usual elimination constructs those representatives. An independent enumeration for \(1\leq M,N\leq4\) confirms the formula and uniqueness of their rank profiles; examples are \(R_{1,1}=2\), \(R_{2,2}=7\), \(R_{2,3}=13\), \(R_{3,3}=34\), and \(R_{4,4}=209\). These finite checks support, but do not replace, the geometric argument. No assertion about the other side of Proposition 4.1.1 or the full maze correspondence is made.\par
\section*{Verification and limits}
An independently written exact polynomial script verifies the Hamiltonian derivatives, the three Lie–Poisson identities, and the nilpotent identity \(\mu_H^2+4\mu_E\mu_F=0\), as well as the finite rook enumeration. The dimension obstruction is exact and analytic, not a numerical conjecture. No author code was executed.\par
The actual source statement and its proof page, the linear-case source pages, and every final report page are displayed and inspected, with exact source, TeX, PDF and compilation bindings in the supporting records. A bounded publisher/arXiv correction search did not identify a correction. Publisher HTML retains the locally closed hypothesis. A preprint retrieval exposed differently organized text; no preprint/PDF equivalence is asserted. Acquisition supplement/code search limitations are retained.\par
Target | Evidence | Impact | Status\par
Problem 1 | Hamiltonian plane, \(Z=\{0\}\) | False equivalence in 2.1.1(1) | Major; initial finding\par
Open-\(Z\) restriction | Obstruction removed in example | Candidate scope correction | Full repair unverified\par
Linear cotangent side | Conormals and rook ranks | Checked independently | No issue found in scope\par
Other component counts | Partial contextual reading | No refutation asserted | Not certified\par
Scientific decision: pending. Human reviews: none. Initial review does not constitute independent acceptance.\par
9 M. Finkelberg, V. Ginzburg and R. Travkin, Lagrangian Subvarieties of Hyperspherical Varieties, Geom. Funct. Anal. 35 (2025), 254–282. https://doi.org/10.1007/s00039-025-00703-3. Assigned version-of-record PDF.\par
\end{document}
